\section{The generic bar-joint rigidity matroid in dimension three}
\label{sec:bar-joint-3d}

The molecular conjecture's application to molecular structures ---
KT Corollary~5.7, the rank formula
$r(G^2) = 3|V| - 6 - \operatorname{def}(\tilde G)$ for the square of a
graph of minimum degree at least two, established by
Jackson--Jord\'an~\cite{jacksonJordan2008} conditional on the
conjecture and unconditionally by
Katoh--Tanigawa~\cite{katohTanigawa2011} --- is a statement about the
\emph{generic bar-joint rigidity matroid} in dimension three. This
chapter constructs that matroid in linear-matroid form, on the
framework substrate of \cref{sec:frameworks} and by the
\texttt{Matroid.ofFun} device of
\cref{sec:rigidity-matroid-linear-framing}.

Only the linear matroid and its rank function are needed in the
sequel. A combinatorial characterization of generic bar-joint rigidity
in dimension three --- the analogue of Laman's theorem
(\cref{sec:laman-theorem}) --- is a well-known open
problem~\cite[\S7]{katohTanigawa2011}, and none is attempted here.

In dimension two, the linear framing of
\cref{sec:rigidity-matroid-linear-framing} was measured against the
combinatorial $(2,3)$-count matroid, and its generic-placement lemma
(\cref{lem:exists-uniform-rowIndependent-placement}) obtained the
per-subset witness placements from the sparsity characterization
(\cref{thm:isSparse-exists-rowIndependent-placement}). In general
dimension no combinatorial description is available, and none is
needed: a placement is \emph{generic for row independence}
(\cref{def:generic-placement}) when every edge set that is
row-independent at some placement is row-independent at it, and such
placements exist in every dimension
(\cref{lem:exists-generic-placement}) by the same linear-interpolation
argument, the witness placements now supplied by the definition of the
family itself. The matroid $\mathcal{R}_d(V)$ at such a placement
(\cref{def:genericRigidityMatroid}) does not depend on the placement
chosen (\cref{lem:linearRigidityMatroid-eq-genericRigidityMatroid}),
its independent sets are exactly the edge sets row-independent at some
placement (\cref{lem:genericRigidityMatroid-indep-iff}), and in
dimension two it recovers the planar rigidity matroid
(\cref{lem:genericRigidityMatroid-two-eq-rigidityMatroid}). The
chapter closes with the generic rank function $r_d$
(\cref{def:genericRank}) and its row-space form
(\cref{lem:genericRank-eq-finrank-span}), the two interfaces the
molecule application reads at $d = 3$.

Throughout, $V$ is a finite vertex set and $d$ an arbitrary dimension;
the molecule application takes $d = 3$, where
$d\,|V| - \binom{d+1}{2} = 3|V| - 6$ is the maximal possible rank of a
spanning graph.

\subsection{Generic placements}
\label{sec:bar-joint-3d-generic-placement}

A placement with algebraically independent coordinates realizes every
row rank that any placement realizes; the definition below asks only
for that consequence, which is all the matroid construction uses, so
existence can be proved by an elementary perturbation argument rather
than by transcendence degree.

\begin{definition}[Generic placement for row independence]
  \label{def:generic-placement}
  \lean{SimpleGraph.IsGenericPlacement}
  \leanok
  \uses{def:framework, def:edgeSet-rowIndependent}
  A placement $p : V \to \R^d$ is \emph{generic for row independence}
  if every edge subset $I \subseteq E(K_V)$ that is row-independent at
  some placement (\cref{def:edgeSet-rowIndependent}, at $G = K_V$) is
  row-independent at $p$.
\end{definition}

\begin{lemma}[Existence of generic placements]
  \label{lem:exists-generic-placement}
  \lean{SimpleGraph.exists_isGenericPlacement}
  \leanok
  \uses{def:generic-placement}
  For every finite $V$ and every dimension $d$ there exists a
  placement $p : V \to \R^d$ generic for row independence.
\end{lemma}
\begin{proof}
  \leanok
  \uses{def:edgeSet-rowIndependent}
  Induction on the cardinality of the finite family
  \[
    \mathcal{F} := \{\, I \subseteq E(K_V) \mid
      I\ \text{is row-independent at some placement}\,\},
  \]
  as in \cref{lem:exists-uniform-rowIndependent-placement}, whose
  dimension-2 statement differs only in producing the witness
  placements from the sparsity characterization: here a member
  $I \in \mathcal{F}$ carries its witness placement by definition.
  Given $p_0$ row-independent on every member of a subfamily and a new
  member $I_0$ with witness $q$, each condition ``$S$ is
  row-independent at $p_t := p_0 + t\,(q - p_0)$'' is the
  non-vanishing of a one-variable polynomial in $t$ (the Gram
  determinant of the rigidity-row matrix, whose rows are affine in
  $t$), nonzero at $t = 0$ for the subfamily members and at $t = 1$
  for $I_0$; a $t$ avoiding the finitely many roots enlarges the
  subfamily.
\end{proof}

\subsection{The matroid and its rank function}
\label{sec:bar-joint-3d-matroid}

\begin{definition}[The generic rigidity matroid $\mathcal{R}_d(V)$]
  \label{def:genericRigidityMatroid}
  \lean{SimpleGraph.genericRigidityMatroid}
  \leanok
  \uses{def:linearRigidityMatroid, lem:exists-generic-placement}
  Fix, by \cref{lem:exists-generic-placement}, a placement
  $p : V \to \R^d$ generic for row independence. The \emph{generic
  ($d$-dimensional bar-joint) rigidity matroid} of $V$ is
  $\mathcal{R}_d(V) := \mathcal{L}_p(V)$, the linear rigidity matroid
  at $p$ (\cref{def:linearRigidityMatroid}) --- a matroid on
  $\mathrm{Sym}_2 V$ with ground set $E(K_V)$. By
  \cref{lem:linearRigidityMatroid-eq-genericRigidityMatroid} the
  matroid does not depend on the choice of $p$.
\end{definition}

\begin{lemma}[Independence in $\mathcal{R}_d(V)$]
  \label{lem:genericRigidityMatroid-indep-iff}
  \lean{SimpleGraph.genericRigidityMatroid_indep_iff}
  \leanok
  \uses{def:genericRigidityMatroid, def:edgeSet-rowIndependent}
  An edge subset $I \subseteq E(K_V)$ is independent in
  $\mathcal{R}_d(V)$ if and only if $I$ is row-independent at some
  placement $q : V \to \R^d$.
\end{lemma}
\begin{proof}
  \leanok
  \uses{lem:linearRigidityMatroid-indep-iff, def:generic-placement}
  Let $p$ be the generic placement defining $\mathcal{R}_d(V)$.
  Independence in $\mathcal{L}_p(V)$ is row independence at $p$
  (\cref{lem:linearRigidityMatroid-indep-iff}). Row independence at
  $p$ exhibits the witness $q = p$; conversely, row independence at
  any witness placement transfers to $p$ by genericity
  (\cref{def:generic-placement}).
\end{proof}

\begin{lemma}[Placement independence]
  \label{lem:linearRigidityMatroid-eq-genericRigidityMatroid}
  \lean{SimpleGraph.linearRigidityMatroid_eq_genericRigidityMatroid}
  \leanok
  \uses{def:genericRigidityMatroid, def:generic-placement}
  For every placement $p : V \to \R^d$ generic for row independence,
  $\mathcal{L}_p(V) = \mathcal{R}_d(V)$.
\end{lemma}
\begin{proof}
  \leanok
  \uses{lem:linearRigidityMatroid-indep-iff,
    lem:genericRigidityMatroid-indep-iff}
  Both matroids have ground set $E(K_V)$, so by
  \texttt{Matroid.ext\_indep} it suffices to identify their
  independent sets. Independence in $\mathcal{L}_p(V)$ is row
  independence at $p$ (\cref{lem:linearRigidityMatroid-indep-iff}),
  which by genericity of $p$ is equivalent to row independence at some
  placement, i.e.\ to independence in $\mathcal{R}_d(V)$
  (\cref{lem:genericRigidityMatroid-indep-iff}).
\end{proof}

\begin{lemma}[Dimension two: the planar rigidity matroid]
  \label{lem:genericRigidityMatroid-two-eq-rigidityMatroid}
  \uses{def:genericRigidityMatroid, def:rigidityMatroid}
  $\mathcal{R}_2(V) = \mathcal{R}(V)$: in dimension two the generic
  rigidity matroid coincides with the combinatorial planar rigidity
  matroid of \cref{def:rigidityMatroid}.
\end{lemma}
\begin{proof}
  \uses{lem:genericRigidityMatroid-indep-iff,
    thm:rigidityMatroid-indep-iff-rowIndependent}
  Both matroids have ground set $E(K_V)$, and both independence
  predicates say ``row-independent at some placement''
  (\cref{lem:genericRigidityMatroid-indep-iff};
  \cref{thm:rigidityMatroid-indep-iff-rowIndependent}).
\end{proof}

\begin{definition}[Generic rank]
  \label{def:genericRank}
  \uses{def:genericRigidityMatroid}
  For a graph $H$ on $V$, the \emph{generic ($d$-dimensional) rank} of
  $H$ is $r_d(H)$, the rank of the edge set $E(H)$ in
  $\mathcal{R}_d(V)$. At $d = 3$ this is the rank function $r(\cdot)$
  of the molecule rank formula (KT Corollary~5.7).
\end{definition}

\begin{lemma}[Row-space form of the generic rank]
  \label{lem:genericRank-eq-finrank-span}
  \uses{def:genericRank, def:generic-placement, def:rigidityRow}
  Let $p : V \to \R^d$ be a placement generic for row independence.
  For every graph $H$ on $V$,
  \[
    r_d(H) \;=\; \dim \operatorname{span}
      \{\,\mathrm{rigidityRow}\,K_V\,p\,e \mid e \in E(H)\,\},
  \]
  the dimension of the span of the rigidity rows of $H$'s edges at
  $p$.
\end{lemma}
\begin{proof}
  \uses{lem:linearRigidityMatroid-eq-genericRigidityMatroid,
    lem:linearRigidityMatroid-indep-iff}
  By \cref{lem:linearRigidityMatroid-eq-genericRigidityMatroid},
  $\mathcal{R}_d(V) = \mathcal{L}_p(V)$, so the rank of $E(H)$ is the
  maximal cardinality of a subset of $E(H)$ on which the row function
  at $p$ is linearly independent
  (\cref{lem:linearRigidityMatroid-indep-iff}). A maximal such subset
  indexes a basis of the row span, so its cardinality is the
  dimension of the span.
\end{proof}
